\chapter{Executive overview}\label{ch:exec}

\section{What CETP is}
CETP (Cross-Environment Runtime Prediction) answers one question for a PostgreSQL user: \emph{a query
ran in $X$\,ms on my development machine; how long will it take on production?} The input is the
\texttt{EXPLAIN (ANALYZE, FORMAT JSON)} of the query on the development machine, a numeric description
of both machines, and (optionally) a handful of timed calibration queries on production. The output
is a scaling factor $s=T_{\text{prod}}/T_{\text{dev}}$, a predicted runtime, and a trust level.
Prediction itself needs no database access. The only step that writes to a database is the optional
capability probe, which creates and then drops three \texttt{UNLOGGED} tables named \texttt{\_cetp\_probe\_*}
(\cref{ch:impl}). The package has no runtime dependencies (\src{pyproject.toml}, \evFILE).

\section{The honest headline}
The repository contains two generations of the method and one cannot summarise CETP without saying
what each does and where each fails. The following statements are the ones I would defend to a panel.

\begin{finding}
\textbf{1. On the original five-machine EC2 study, a two-scalar analytical model beats the
do-nothing baseline by a small, statistically supported margin.} Leave-one-machine-out MAPE: naive
31.96\%, naive-linear 23.05\%, analytical 21.61\%. The paired difference naive-linear minus analytical
is $+1.44$\,pp with 95\% cluster-bootstrap CI $[0.32,\,2.56]$ \evFILE\
(\src{results/canonical_metrics.json}). The analytical model is \emph{not} distinguishable from the
even simpler compute-only rule: difference $+1.55$\,pp, CI $[-0.41,\,3.51]$.
\end{finding}

\begin{finding}
\textbf{2. A per-operator model that is aware of PostgreSQL parallelism wins decisively over
scalar scaling where parallelism changes between machines.} On the regenerated Docker grid (8
arm64 cells, 22 queries, 1{,}232 pairs) MAPE falls from 66.97\% (scalar, geometric fallback) to
30.15\% once the \emph{parallel} scaling is added and to 22.41\% for the full per-operator model; the
naive baseline is 72.52\% \evFILE\ (\src{results/operator_model_metrics.json}). The mechanism is that
a parallel scan sees a speed-up of about $2.7\times$ when the CPU quota allows three workers, and about
$0.77\times$ when it does not.
\end{finding}

\begin{warnbox}
\textbf{3. Zero-shot cross-ISA prediction currently fails.} Predicting x86\_64 GitHub-Actions runtimes
from an arm64 laptop, the operator model has MAPE 66--73\% against 21--23\% for the naive baseline
(mean 68.92\% vs 21.74\%) \evFILE\ (\src{results/laptop_to_prod_metrics.json}). Over all 11 targets
the operator model is worse than doing nothing (43.72\% vs 31.70\%); only on the 7 same-ISA targets is it
better (29.31\% vs 37.39\%). The library function therefore \emph{refuses} to predict across architectures without
calibration (\texttt{CrossArchitectureUncalibrated}, \evCODE\ \src{cetp/cross\_env.py}); but the
\texttt{cetp predict-prod} command does not pass the information the refusal needs and returns a
prediction anyway (\evRUN, \cref{sec:cli-bypass}).
\end{warnbox}

\begin{warnbox}
\textbf{4. Laptop-to-EC2 transfer does not work for either method.} On 95 pairs, naive error is 129.39\%
and the cheaper laptop-derived signatures are far worse (160.68\% and 337.75\%)
\evFILE\ (\src{results/laptop_to_ec2_metrics.json}). The operator model cannot be evaluated here because the
EC2 hosts were terminated.
\end{warnbox}

\section{What is genuinely contributed}
The repository's own research pack recommends (and I agree) that the defensible contribution is
\emph{not} either named method. It is:
\begin{enumerate}
\item the finding that \textbf{parallel-aware scaling is the capability that matters}: restricting the
  hardware term to plan nodes beneath a \texttt{Gather} reduces error from 66.97\% to 22.23--22.41\%;
\item the finding that \textbf{cross-ISA transfer fails and must be gated} by calibration or by refusal;
\item a set of \textbf{negative results reported with confidence intervals}: the roofline weighting and
  the per-operator weighting are statistically indistinguishable from a single scalar once parallelism is
  handled ($-0.17$\,pp, CI $[-0.49,\,+0.10]$), and a learned model does not beat the analytical formula.
\end{enumerate}
It is designed to be reproducible: canonical numbers are regenerated by one script and the Makefile compares two runs byte-for-byte (the byte-identical outcome is \evPACK; I did not re-run it).

\section{What it is not}
\begin{itemize}
\item It is not a production-ready predictor: HEAD is an alpha package, with no measured end-to-end
  accuracy on a real production database. The only ``production'' evidence is four GitHub-Actions runners
  and eight local containers.
\item It does not predict the plan: it assumes the production planner picks the same plan. On the Docker
  grid 0 of 22 plans differ across cells, so plan changes are untested there; across PostgreSQL versions
  12 of 22 plans change (\evFILE\ \src{results/tier2\_versions.csv}, as reported by the research pack).
\item It does not use hardware performance counters; AWS Nitro hides them (\cref{ch:problem}).
\end{itemize}

\section{System at a glance}
\Cref{fig:overview} shows the end-to-end flow; the numbers in each box are the sections that explain it.

\begin{figure}[ht]
\centering
\begin{tikzpicture}[node distance=6mm and 9mm,
  box/.style={draw=cetpblue,rounded corners=2pt,fill=lightgrey,minimum height=9mm,text width=26mm,align=center,font=\footnotesize},
  io/.style={draw=cetpteal,rounded corners=2pt,fill=cetpteal!8,minimum height=9mm,text width=26mm,align=center,font=\footnotesize},
  arr/.style={-{Stealth[length=2mm]},thick,cetpgrey}]
\node[io] (dev) {Dev machine\\ \texttt{EXPLAIN ANALYZE}};
\node[box,right=of dev] (dec) {Decompose into per-operator self-times};
\node[box,right=of dec] (w) {Weights and \texttt{Gather} flag};
\node[box,right=of w] (sig) {Price with machine rates (probe)};
\node[io,below=of sig] (pred) {Scaling $s$, runtime, trust};
\node[box,left=of pred] (cal) {Calibration (median actual / predicted)};
\node[io,left=of cal] (few) {$k$ timed queries on production};
\draw[arr] (dev)--(dec); \draw[arr] (dec)--(w); \draw[arr] (w)--(sig); \draw[arr] (sig)--(pred);
\draw[arr] (few)--(cal); \draw[arr] (cal)--(pred);
\end{tikzpicture}
\caption{Gen-2 (per-operator) prediction path. Derived from \src{cetp/plan\_nodes.py}, \src{cetp/operator\_model.py} and \src{cetp/cross\_env.py}.}
\label{fig:overview}
\end{figure}

\section{Two generations, one repository}\label{sec:two-gens}
\begin{center}\small
\begin{tabularx}{\textwidth}{@{}L{2.4cm}YY@{}}
\toprule \thead{} & \thead{Gen 1 ``roofline''} & \thead{Gen 2 ``per-operator''}\\ \midrule
Machine signature & 2 scalars: STREAM bandwidth, multi-core GFLOPS & Per-operator rates from a PostgreSQL micro-probe; parallel speed-up\\
Query model & Whole-query: one global factor & Decomposed plan: weights over operators, \texttt{Gather} aware\\
Formula & $\sigma = 0.5/B + 0.25/C + 0.25\sqrt{1/(BC)}$ & $\sigma=\sum_b w_b\,\rho_b\,\pi_b\,\kappa$ (\cref{ch:algo})\\
Evidence & 5 EC2 machines, 840 pairs & 8 Docker cells (1{,}232 pairs); 11 laptop-to-prod cells (242 pairs)\\
Exposed in & CLI, VS Code extension, CI gate & CLI (\texttt{predict}, \texttt{calibrate}, \texttt{check}); not the extension or CI gate\\
\bottomrule
\end{tabularx}
\end{center}
The Gen-1 bandwidth $B$ and compute $C$ are ratios $\text{prod}/\text{dev}$. The ``roofline'' label is
historical: the formula is a fixed-weight blend, not the minimum of a compute and a bandwidth bound
\citep{williams:roofline}. A reader who finds ``roofline'' in the code or older decks should read it as
``Gen-1 analytical''.

\section{Reading guide to the numbers}
Two result sets are quoted in this repository: the \emph{canonical} numbers
(\src{results/canonical_metrics.json}, regenerated byte-identically by one script) and older ``corrected''
numbers (\src{results/corrected_headline_mape.json}). They differ in the second decimal
(analytical 21.61 vs 21.599; naive-linear 23.05 vs 23.025). This document uses canonical values and
labels older ones. The Review 3 slide deck uses the older values; this is not an error of substance but
the two must not be mixed in one table (\cref{app:issues}).
